\documentclass[12pt]{article}
\usepackage[utf8]{inputenc}
\usepackage[T1]{fontenc}
\usepackage[margin=1in]{geometry}
\usepackage{setspace}
\usepackage{xcolor}
\usepackage[colorlinks=true, urlcolor=blue!60!black, linkcolor=blue!60!black]{hyperref}
\setlength{\parindent}{0em}
\setlength{\parskip}{0.45em}
\renewcommand{\baselinestretch}{1.0}
\begin{document}
\begin{center}
{\Large \textit{Heart of Atlanta Motel, Inc.\ v.\ United States}, 379 U.S.\ 241 (1964)\par}
\vspace{0.3em}
{\normalsize Case Brief --- One-Page Edition\par}
\end{center}

\textbf{Memory Jogger.} Congress may use the Commerce Clause to ban racial discrimination by motels serving interstate travelers (\href{https://supreme.justia.com/cases/federal/us/379/241/}{Justia}).

\textbf{Facts.} Title II of the Civil Rights Act of 1964 barred discrimination in public accommodations affecting commerce. The Heart of Atlanta Motel, near two interstates and drawing about 75\% of its guests from out of state, refused Black guests and challenged the Act.

\textbf{Procedural History.} A three-judge district court upheld Title II and enjoined the motel; direct appeal.

\textbf{Votes.} Affirmed, 9--0: Clark for the Court; Black, Douglas, and Goldberg concurred.

\textbf{Holding.} Congress may bar discrimination by public accommodations serving interstate travelers when it rationally finds an effect on commerce, even while remedying a moral wrong.

\textbf{Reasoning.} The \textit{Civil Rights Cases} did not control. Hearings showed that discrimination deterred Black interstate travel; the movement of persons is commerce, and Congress may reach its local obstructions. The Fifth and Thirteenth Amendment claims failed. Black relied on aggregation; Douglas and Goldberg would also have rested the Act on \S~5 of the Fourteenth Amendment.

\textbf{Examples.} \emph{Same-Side:} commercial vacation-rental hosts. \emph{Opposite-Side:} local gender-motivated violence (\textit{Morrison}). \emph{Fence-Sitter:} a remote farmhouse with in-state guests.

\textbf{Critique.} Originalists see a police power; progressives regret the commercial framing.

\textbf{Key Quotation.} ``Thus, the power of Congress to promote interstate commerce also includes the power to regulate the local incidents thereof, including local activities in both the States of origin and destination, which might have a substantial and harmful effect upon that commerce.'' --- Justice Clark, 379 U.S.\ at 258.

\end{document}
